\documentclass[aspectratio=169, 11pt]{beamer}

% --- Theme & Colors ---
\usetheme{Madrid}
\usecolortheme{seahorse}
\setbeamertemplate{navigation symbols}{}
\setbeamertemplate{footline}[frame number]

% --- Packages ---
\usepackage{amsmath, amssymb, amsthm}
\usepackage{graphicx}
\usepackage{tikz}
\usepackage{booktabs}

% --- Custom colors ---
\definecolor{tseblue}{RGB}{0,70,130}
\setbeamercolor{title}{fg=tseblue}
\setbeamercolor{frametitle}{fg=tseblue}
\setbeamercolor{structure}{fg=tseblue}

% --- Title ---
\title[Armstrong \& Zhou (2022)]{Consumer Information and the Limits to Competition}
\subtitle{Armstrong \& Zhou, \textit{American Economic Review}, 2022}
\author{Jordi Torres Vallverd\'{u}}
\institute{Toulouse School of Economics}
\date{IO Theory -- MRes 2025--2026}

\begin{document}

% =============================================
% TITLE
% =============================================
\begin{frame}
\titlepage
\end{frame}

% =============================================
% PART 1: CONTEXT / PROBLEM (~5 min)
% =============================================

% --- Slide: Why This Matters ---
\begin{frame}{Why Does This Matter?}

\textbf{Platforms design the information environment consumers face}

\bigskip

\begin{itemize}
    \item Amazon, Expedia, insurance comparison sites choose:
    \begin{itemize}
        \item How much product detail to reveal
        \item Whether to offer personalized recommendations
        \item How flexibly consumers can filter and compare
    \end{itemize}

    \bigskip

    \item \textbf{Core question:} Should a platform (or regulator) make products \emph{more} or \emph{less} comparable?

    \bigskip

    \item More transparency $\Rightarrow$ better product match, but firms gain market power $\Rightarrow$ higher prices

    \item Less transparency $\Rightarrow$ fierce price competition, but consumers may buy the wrong product
\end{itemize}

\end{frame}

% --- Slide: Main Idea & Tension ---
\begin{frame}{Main Idea and Key Tension}

\textbf{Central trade-off:} \; match quality \; vs.\ price competition

\bigskip

\begin{columns}[T]
\column{0.48\textwidth}
\textbf{Firm-optimal policy}
\begin{itemize}
    \item Amplifies perceived differentiation
    \item Fewer price-sensitive consumers
    \item Higher equilibrium prices
    \item \alert{No mismatch} $\Rightarrow$ max total welfare
\end{itemize}

\column{0.48\textwidth}
\textbf{Consumer-optimal policy}
\begin{itemize}
    \item Dampens perceived differentiation
    \item More price-sensitive consumers
    \item Lower equilibrium prices
    \item \alert{Some mismatch} $\Rightarrow$ welfare loss
\end{itemize}
\end{columns}

\bigskip

\centering
\textit{Firms and consumers prefer opposite signal structures --- \\
but both prefer symmetric signals over biased recommendations}

\end{frame}

% --- Slide: Literature ---
\begin{frame}{Contribution to the Literature}

\textbf{Information design meets oligopoly pricing}

\bigskip

\begin{itemize}
    \item \textbf{Monopoly information design} (Roesler \& Szentes, 2017)
    \begin{itemize}
        \item[$\rightarrow$] Extends to duopoly; competition changes the problem fundamentally
    \end{itemize}

    \medskip

    \item \textbf{Bayesian persuasion} (Kamenica \& Gentzkow, 2011)
    \begin{itemize}
        \item[$\rightarrow$] Applies persuasion framework to competitive pricing
    \end{itemize}

    \medskip

    \item \textbf{Competitive disclosure} (Anderson \& Renault, 2006; Johnson \& Myatt, 2006)
    \begin{itemize}
        \item[$\rightarrow$] Goes beyond ``truth or noise'' to general signal structures
    \end{itemize}

    \medskip

    \item \textbf{Centralized vs.\ decentralized design} (Ivanov, 2013; Hwang et al., 2019)
    \begin{itemize}
        \item[$\rightarrow$] Centralized design allows \emph{relative} signals (e.g., ``top product'')
    \end{itemize}
\end{itemize}

\end{frame}

% =============================================
% PART 2: MODEL (~5 min)
% =============================================

% --- Slide: Setup ---
\begin{frame}{Model Setup}

\begin{itemize}
    \item Two firms ($i = 1,2$) supply differentiated products at zero cost

    \medskip

    \item Consumer values product $i$ at $v_i \geq 0$; prior on $(v_1, v_2)$ is symmetric

    \medskip

    \item Only the \textbf{relative preference} matters: $x \equiv v_1 - v_2$
    \begin{itemize}
        \item Prior CDF $F(x)$ symmetric around 0, support $[-1,1]$
    \end{itemize}

    \medskip

    \item Consumer observes a \textbf{private signal} $s$ before purchase
    \begin{itemize}
        \item Signal structure $\{\sigma(s|x), S\}$ is public and chosen \emph{before} firms set prices
    \end{itemize}

    \medskip

    \item After signal $s$, consumer updates to \textbf{posterior} $G(x) \equiv$ CDF of $E_{F_s}[x]$

    \medskip

    \item Firms observe $G$ but not $s$; set prices simultaneously (Bertrand--Nash)

    \medskip

    \item Consumer buys from firm 1 if $E[x \mid s] > p_1 - p_2$
\end{itemize}

\end{frame}

% --- Slide: Key Constraint ---
\begin{frame}{Key Constraint: Mean-Preserving Contraction}

\textbf{Bayesian consistency:} any posterior $G$ must be a \textbf{mean-preserving contraction} (MPC) of the prior $F$:

$$\int_{-1}^{x} G(\tilde{x})\, d\tilde{x} \;\leq\; \int_{-1}^{x} F(\tilde{x})\, d\tilde{x} \qquad \text{for all } x \in [-1,1]$$

\bigskip

\begin{itemize}
    \item Any $G$ that is an MPC of $F$ can be induced by \emph{some} signal structure

    \medskip

    \item This allows us to work directly with the posterior $G$ rather than the signal

    \medskip

    \item Key parameter: $\delta \equiv E_F[\max\{x, 0\}]$ -- the \textbf{match efficiency under full info}
    \begin{itemize}
        \item Captures the degree of underlying product differentiation
    \end{itemize}
\end{itemize}

\end{frame}

% --- Slide: Bounds ---
\begin{frame}{Bounds on Equilibrium Posteriors (Lemma 1)}

\begin{block}{Lemma 1}
A distribution $G$ implements symmetric prices $(p_1 = p_2 = p)$ if and only if
$$L_p(x) \;\leq\; G(x) \;\leq\; U_p(x) \qquad \text{for all } x \in [-1,1]$$
\end{block}

\medskip

where the bounds arise from each firm's \textbf{no-deviation condition}:

$$L_p(x) = \begin{cases} 0 & \text{if } x \leq -\tfrac{1}{2}p \\ 1 - \dfrac{p}{2(p+x)} & \text{otherwise} \end{cases}$$

$$U_p(x) = \begin{cases} \dfrac{p}{2(p-x)} & \text{if } x < \tfrac{1}{2}p \\ 1 & \text{otherwise} \end{cases}$$

\medskip

\small
Equilibrium profit: $\pi_i = p_i^2/(p_1 + p_2)$. \; Market shares and profits depend only on prices, not on $G$.

\end{frame}

% --- Slide: Symmetry Result ---
\begin{frame}{Symmetric Signals Dominate (Lemma 2)}

\begin{block}{Lemma 2}
Relative to any \emph{asymmetric} posterior (one that treats firms differently):
\begin{enumerate}
    \item[(i)] There exists a \textbf{symmetric} posterior with \textbf{higher profit for each firm}
    \item[(ii)] There exists a \textbf{symmetric} posterior with \textbf{higher consumer surplus}
\end{enumerate}
\end{block}

\bigskip

\textbf{Intuition:}
\begin{itemize}
    \item Biased signals favor one firm $\Rightarrow$ the other undercuts aggressively
    \item The resulting price war hurts \emph{both} firms --- even the favored one
    \item For consumers: asymmetry causes mismatch + price dispersion
\end{itemize}

\bigskip

$\Rightarrow$ \textbf{We can restrict attention to symmetric signal structures}

\end{frame}

% --- Slide: Proposition 1 ---
\begin{frame}{Firm-Optimal Signal Structure (Proposition 1)}

\begin{block}{Proposition 1}
Under Condition 2(i), the firm-optimal signal structure induces the symmetric equilibrium price:
$$p^* \;=\; \frac{2\delta}{1 - \log 2} \;\approx\; 6.5\,\delta$$
It is uniquely implemented by $G(x) = L_{p^*}(x)$ for $x < 0$. \\[4pt]
\textbf{No mismatch occurs} and \textbf{total welfare is maximized}.
\end{block}

\bigskip

\textbf{Key properties:}
\begin{itemize}
    \item The posterior ``pushes'' consumers toward the extremes (U-shaped density)
    \item Fewer marginal consumers $\Rightarrow$ relaxed price competition
    \item Consumer always buys her preferred product
    \item Profit at least 63\% higher than under full disclosure (Corollary 1)
\end{itemize}

\end{frame}

% --- Slide: Proposition 2 ---
\begin{frame}{Consumer-Optimal Signal Structure (Proposition 2)}

\begin{block}{Proposition 2}
Under Conditions 1 and 2(ii), the consumer-optimal signal induces:
$$p^{**} \;=\; \frac{-\gamma}{1 - \gamma}\, F^{-1}\!\Big(\frac{\gamma}{2}\Big), \qquad \gamma \approx 0.05$$
where $\gamma$ solves $\gamma - \log \gamma = 3$. \\[4pt]
Implemented by $G(x) = \min\{F(x),\, U_{p^{**}}(x)\}$ for $x < 0$. \\[4pt]
\textbf{There is mismatch} and \textbf{total welfare is not maximized}.
\end{block}

\bigskip

\textbf{Key properties:}
\begin{itemize}
    \item The posterior amplifies the mass of \emph{marginal} consumers (near $x = 0$)
    \item Firms forced to compete on price $\Rightarrow$ very low equilibrium price
    \item Consumers with weak preferences may buy the wrong product
    \item Price effect dominates the match quality effect for consumers
\end{itemize}

\end{frame}

% =============================================
% PART 3: DISCUSSION (~5 min)
% =============================================

% --- Slide: Feasible Frontier ---
\begin{frame}{The Limits to Competition (Proposition 3)}

\begin{block}{Proposition 3}
Any profit $p \in [0, p^*]$ is feasible. For a given profit $p \leq p^*$:
\begin{itemize}
    \item \textbf{Min consumer surplus:} $\mu - \tfrac{1}{2}(1 + \log 2)\, p$
    \item \textbf{Max consumer surplus:} $\mu + \delta - p$ \; (when $p \geq p_F$), concave in $p$ when $p < p_F$
\end{itemize}
\end{block}

\bigskip

\textbf{Implications:}
\begin{itemize}
    \item Full disclosure is optimal \emph{only} when maximizing unweighted total welfare
    \item Any weighted sum (Ramsey problem) picks a point on the feasible frontier
    \item Due to competition, it is \textbf{impossible} to have both low profit \emph{and} low consumer surplus (contrast with monopoly)
\end{itemize}

\end{frame}

% --- Slide: Extensions ---
\begin{frame}{Extensions}

\textbf{1. Mixed-strategy equilibria} (Section IIIA)
\begin{itemize}
    \item Allowing mixed strategies improves consumer surplus by at most 1.6\%
    \item Pure-strategy restriction is essentially without loss
\end{itemize}

\bigskip

\textbf{2. More than two firms} (Section IIIB)
\begin{itemize}
    \item ``Top product'' signal: tell consumer her best match only
    \begin{itemize}
        \item Can achieve first-best profit; dominates full disclosure for firms
    \end{itemize}
    \item ``Top two'' signal: tell consumer her best \emph{two} matches (no ranking)
    \begin{itemize}
        \item Induces Bertrand competition $\Rightarrow$ price = 0
        \item Dominates both ``top product'' and full disclosure for consumers
    \end{itemize}
    \item As $n \to \infty$: trade-off between price and match quality vanishes
\end{itemize}

\end{frame}

% --- Slide: Criticism & Future ---
\begin{frame}{Discussion and Open Questions}

\textbf{Strengths}
\begin{itemize}
    \item Clean characterization of the fundamental trade-off
    \item Working with posteriors (via MPC) makes the problem tractable
    \item Strong dominance of symmetric signals: surprising and policy-relevant
\end{itemize}

\bigskip

\textbf{Limitations and directions}
\begin{itemize}
    \item Scalar heterogeneity ($x = v_1 - v_2$): multi-dimensional preferences are harder
    \item Symmetric firms: what if one firm is known to be better?
    \item Information designer as strategic agent (platform with own objective)?
    \item Endogenous product design: firms may adjust products in response to signals
    \item Empirical testing: how do real platforms' information policies compare?
\end{itemize}

\end{frame}

% --- Slide: Takeaways ---
\begin{frame}{Key Takeaways}

\begin{enumerate}
    \item \textbf{More information $\neq$ more competition}
    \begin{itemize}
        \item Firm-optimal policy gives consumers \emph{more} info than full disclosure (in the sense of higher prices) while ensuring perfect matching
    \end{itemize}

    \bigskip

    \item \textbf{Consumers prefer less precise information}
    \begin{itemize}
        \item Lower prices outweigh the cost of occasional mismatch
    \end{itemize}

    \bigskip

    \item \textbf{Biased recommendations never help}
    \begin{itemize}
        \item Symmetric signals dominate for firms \emph{and} consumers
    \end{itemize}

    \bigskip

    \item \textbf{Information design has limits}
    \begin{itemize}
        \item Competition prevents simultaneously low profit and low consumer surplus
    \end{itemize}
\end{enumerate}

\end{frame}

% --- Thank you ---
\begin{frame}
\centering
\vfill
{\Large \textbf{Thank you}}

\bigskip

{\large Questions?}
\vfill
\end{frame}

\end{document}
